% !PRINTED 17
\weberpage{17}
%<<CONT
\[
 x+(-x)=0,
 \qquad x+0=x,
\]
\[
 x+(-y)=x-y,\quad \text{wenn }x>y,
 \quad =-(y-x),\quad \text{wenn }y>x.
\]
Bei dieser Erklärung der Addition gelten, wenn $z_1,z_2,z_3$ irgend
drei Zahlen des ganzen Zahlenbereiches sind, die Gesetze:
\[
 z_1+z_2=z_2+z_1,
 \qquad (z_1+z_2)+z_3=z_1+(z_2+z_3),
\]
die man das commutative und das associative Gesetz nennt. Man bildet
die Summe aus einer beliebigen Anzahl von Summanden, indem man nach
Belieben der Reihe nach je zwei Summanden zu einer Summe vereinigt. Die
Subtraction braucht nicht mehr besonders berücksichtigt zu werden,
wenn man $z_1-z_2$ durch $z_1+(-z_2)$ erklärt und $-(-z)=z$ setzt.

Man stellt die Zahlenreihe anschaulich durch Punkte dar, indem man von
einem festen mit $0$ bezeichneten Punkte einer geraden Linie die
positiven Zahlen als Strecken nach der einen, etwa der rechten, die
negativen Zahlen nach der anderen (linken) Seite aufträgt. Das Bild der
Summe zweier Strecken $z_1+z_2$ erhält man, wenn man von dem Punkte
$z_1$ aus die Strecke von der Länge $\pm z_2$ nach der rechten oder
nach der linken Seite abträgt, je nachdem $z_2$ positiv oder negativ
ist.

Die Multiplication und Division wird in den erweiterten Zahlenbereich
durch die Gleichungen
\[
 x(-y)=(-x)y=-xy,
\]
\[
 (-x)(-y)=xy,
 \qquad 0x=0
\]
erklärt. Die Division als die Umkehrung der Multiplication ist immer
möglich, ausser wenn der Divisor Null ist.

\begin{center}
\rule{0.18\linewidth}{0.4pt}
\end{center}

Eine fernere Erweiterung des Zahlbegriffes besteht in der Einführung
der \gesperrt{complexen Grössen}. Wir combiniren je zwei Zahlen der
gesammten Zahlenreihe zu Paaren $(x,y)$, und betrachten zwei solche
Paare $(x,y)$ und $(a,b)$ nur dann als gleich, wenn $x=a$, $y=b$ ist.
Diese Zahlenpaare bilden eine Mannigfaltigkeit, deren Elemente zwar
nicht geordnet werden, mit denen aber die Rechenoperationen der
Addition, Subtraction, Multiplication und Division vorgenommen werden
sollen, nach folgenden Regeln. Es sei
\[
 (x,y)+(a,b)=(x+a,y+b),
\]
\[
 (x,y)(a,b)=(xa-yb,xb+ya),
\]
%CONT>>
